\documentclass[11pt,a4paper]{article}

\usepackage{amsmath,amssymb,amsthm}
\usepackage{authblk}
\usepackage[margin=2.5cm]{geometry}
\usepackage{hyperref}
\usepackage{booktabs}

\newtheorem{theorem}{Theorem}
\newtheorem{proposition}[theorem]{Proposition}
\newtheorem{corollary}[theorem]{Corollary}
\theoremstyle{definition}
\newtheorem{definition}[theorem]{Definition}
\newtheorem{assumption}[theorem]{Assumption}
\theoremstyle{remark}
\newtheorem{remark}[theorem]{Remark}

\begin{document}

\title{Discrete Mixing Operators from Boundary Sector Geometry}

\author[1]{Marvin L. Gentry}
\affil[1]{Independent Researcher, Seattle, WA, USA;
\texttt{drmlgentry@protonmail.com};
ORCID: 0009-0006-4550-2663}
\date{\today}

\begin{abstract}
The CKM and PMNS flavor mixing matrices of the Standard Model
lack a first-principles geometric origin.
We show that if fermion flavors are associated with distinct boundary
directions of a compact hyperbolic 3-manifold, the resulting mixing
matrices are uniquely determined --- up to sector-local rephasings ---
by the boundary configuration alone, and form a \emph{discrete} set
in $\mathrm{U}(n)/(\mathrm{U}(1)^n\times\mathrm{U}(1)^n)$.
Our main result (Theorem~\ref{thm:discrete}) establishes that no
continuous deformation of physical parameters can produce continuous
drift in any rephasing-invariant mixing observable.
The construction is purely kinematic, requiring no dynamical input,
and provides a model-independent structural constraint on any geometric
theory of flavor mixing.
In companion papers~\cite{Gentry:CKM,Gentry:PMNS}, this framework
is applied to two specific arithmetic hyperbolic 3-manifolds,
reproducing the observed CKM and PMNS mixing matrices with
Frobenius fitness $0.016482$ and $0.005087$ respectively.
\end{abstract}

\maketitle

\section{Introduction}

The CKM quark mixing matrix~\cite{Cabibbo,KM} and the PMNS lepton
mixing matrix~\cite{Maki,Pontecorvo} encode the mismatch between
flavor and mass eigenstates. Despite their precise experimental
determination~\cite{PDG2024}, their values have no derivation from
the gauge structure of the Standard Model. Discrete flavor
symmetries~\cite{AltarelliFeruglio}, extra dimensions, and
string-inspired constructions~\cite{KingLuhn} have all been proposed,
but a kinematic, model-independent explanation has been lacking.

In this paper we take a minimal perspective: fermions are assigned
to distinct boundary directions $\xi(f)\in S^2=\partial\mathbb{H}^3$
of a compact hyperbolic 3-manifold $M=\mathbb{H}^3/\Gamma$.
The mixing matrix between any two sectors arises as the overlap
operator between the corresponding boundary-localized Hilbert
subspaces. We prove that under a discrete assignment rule, the set
of admissible mixing matrices is finite modulo rephasings ---
so continuous drift in any rephasing-invariant observable is
geometrically forbidden.

This result is the kinematic foundation for the Hyperbolic Flavor
Geometry (HFG) program~\cite{Gentry:Unified}.
In the companion CKM paper~\cite{Gentry:CKM}, the K-factor of the
Iwasawa decomposition $\mathrm{PSL}(2,\mathbb{C})=KAN$ applied to
the holonomy of $M_{\mathrm{CKM}}=m006(-5,2)$ gives Frobenius
fitness $0.016482$ against PDG~2024 values.
In the companion PMNS paper~\cite{Gentry:PMNS}, the N-factor (Borel)
applied to $M_{\mathrm{PMNS}}=m003(-2,3)$ gives fitness $0.005087$
--- the global minimum of the construction.
The present paper provides the mathematical guarantee that these
results are not accidents of continuous parameter tuning: the
admissible mixing matrices form a discrete set.

\section{Geometric Setup}

Let $\Gamma$ be a torsion-free discrete group acting properly
discontinuously and cocompactly on $\mathbb{H}^3$ by
orientation-preserving isometries. The quotient $M=\mathbb{H}^3/\Gamma$
is a compact oriented hyperbolic 3-manifold with boundary at
infinity $S^2=\partial\mathbb{H}^3$.

\begin{definition}[Boundary direction]
A \emph{boundary direction} is a point $\xi\in S^2$.
For $\gamma\in\Gamma$, the associated boundary direction is the
ideal endpoint of the geodesic ray from a fixed basepoint $o\in\mathbb{H}^3$
through $\gamma\cdot o$.
\end{definition}

\begin{assumption}[Discrete sector assignment]
\label{ass:discrete}
Each fermion flavor $f\in\mathcal{F}$ is assigned a unique boundary
direction $\xi(f)\in S^2$ by a deterministic, discrete selection rule.
The image $\Xi=\{\xi(f):f\in\mathcal{F}\}$ is finite, and the
assignment is locally rigid: distinct flavors map to distinct
directions, and small perturbations of the basepoint do not change
the assignment class.
\end{assumption}

\begin{remark}[Geometric realization]
A natural realization: for each flavor $f$, let $\gamma_f\in\Gamma$
be the unique element minimizing $d(o,\gamma\cdot o)$ in the
assignment class of $f$. Proper discontinuity of the action implies
the orbit $\Gamma\cdot o$ is discrete, so the minimizer exists;
torsion-freeness forces uniqueness. The boundary direction $\xi(f)$
is then the ideal endpoint of the geodesic from $o$ through
$\gamma_f\cdot o$. In the HFG application, $\gamma_f$ is a short
word in $\pi_1(M)$ and $\xi(f)=\hat{n}(\gamma_f)$ is the
geodesic axis direction extracted via the Pauli decomposition of
$\log\rho(\gamma_f)$.
\end{remark}

\section{Boundary Hilbert Space and Sector Subspaces}

Let $\mathcal{H}$ be a separable complex Hilbert space carrying
boundary data (concretely, $L^2(S^2)$).

\begin{definition}[Boundary-localized state]
For each $\xi\in S^2$, fix a normalized vector $|\xi\rangle\in\mathcal{H}$
such that $\xi\mapsto|\xi\rangle$ is continuous, $\langle\xi|\xi\rangle=1$,
and $|\langle\xi|\zeta\rangle|<1$ for $\xi\neq\zeta$.
An explicit family is the Poisson kernel:
$|\xi\rangle\propto P(\cdot,\xi)$ normalized in $L^2(S^2)$,
which satisfies all three conditions~\cite{Ratcliffe}.
\end{definition}

\begin{definition}[Sector and sector subspace]
A \emph{sector} $S$ is a finite ordered subset of $\mathcal{F}$.
The \emph{sector subspace} is
$\mathcal{H}_S = \mathrm{span}\{|\xi(f)\rangle : f\in S\}\subset\mathcal{H}$,
with $\dim\mathcal{H}_S=|S|$ (since distinct boundary directions give
linearly independent states).
The \emph{preferred ONB} $\{|S,1\rangle,\ldots,|S,n\rangle\}$ of
$\mathcal{H}_S$ is the Gram-Schmidt orthonormalization of
$(|\xi(f_1)\rangle,\ldots,|\xi(f_n)\rangle)$ with all pivot
coefficients real and positive, removing the $\mathrm{U}(1)^n$
phase ambiguity.
\end{definition}

\section{Mixing Matrices as Overlap Operators}

\begin{definition}[Mixing matrix]
For sectors $S$, $T$ of equal cardinality $n$, the
\emph{mixing matrix} $U_{ST}$ is the $n\times n$ matrix with
entries $(U_{ST})_{ji}=\langle T,j|S,i\rangle$.
\end{definition}

\begin{proposition}[Unitarity and geometric dependence]
\label{prop:unitary}
When $\mathcal{H}_S=\mathcal{H}_T$, $U_{ST}\in\mathrm{U}(n)$.
In general it is a partial isometry.
$U_{ST}$ is determined entirely by the finite boundary
configuration $\Xi_{ST}=\{\xi(f):f\in S\cup T\}$ and the inner
product on $\mathcal{H}$: all Gram-Schmidt coefficients are
algebraic expressions in the inner products
$\langle\xi(f)|\xi(g)\rangle$, $f,g\in S\cup T$.
\end{proposition}

\section{Rephasing Invariance}

Physical observables are invariant under sector-local rephasings:
$U_{ST}\to D_T U_{ST} D_S^\dagger$, where $D_S,D_T\in\mathrm{U}(1)^n$.
The physical content is the equivalence class of $U_{ST}$ in
$\mathrm{U}(n)/(\mathrm{U}(1)^n\times\mathrm{U}(1)^n)$.

\begin{proposition}[Rephasing-invariant quantities]
The following are invariant under sector-local rephasings:
(i)~the moduli $|(U_{ST})_{ji}|^2$ (mixing probabilities);
(ii)~the quartic products
$J_{ji,j'i'}=(U_{ST})_{ji}(U_{ST})_{j'i'}(U_{ST})^*_{ji'}(U_{ST})^*_{j'i}$
(generalizing the Jarlskog invariant~\cite{Jarlskog});
(iii)~$\det(U_{ST})$ up to an overall phase.
For $n=3$ and the standard CKM/PMNS conventions, $\mathrm{Im}(J_{ji,j'i'})$
reduces to the standard Jarlskog invariant $J_{\mathrm{CP}}$.
\end{proposition}

\section{Discreteness and Stability}

\begin{theorem}[Stability under deformations]
\label{thm:stable}
Let $\delta_{\min}=\min_{f\neq f'}d_{S^2}(\xi(f),\xi(f'))$ be the
minimal pairwise separation of boundary directions, and set
$\varepsilon_0=\delta_{\min}/2$.
For any deformation of magnitude $\varepsilon\leq\varepsilon_0$
(each $\xi(f)$ moved by at most $\varepsilon$ in the round metric
on $S^2$), the perturbed configuration determines the same preferred
bases and the same mixing matrix $U_{ST}$.
\end{theorem}

\begin{proof}
For $\varepsilon\leq\varepsilon_0$, each perturbed point $\xi'(f)$
remains closer to $\xi(f)$ than to any other $\xi(g)$, $g\neq f$,
so the points remain distinct and carry the same labels.
Linear independence is therefore preserved.
The Gram-Schmidt coefficients depend continuously on the inner
products $\langle\xi'(f)|\xi'(g)\rangle$, which vary continuously
with $\varepsilon$; but the discrete assignment rule
(Assumption~\ref{ass:discrete}) is unchanged since each $\xi'(f)$
belongs to the same assignment class as $\xi(f)$.
By Proposition~\ref{prop:unitary}, $U_{ST}$ depends only on the
configuration class, which is unchanged.
\end{proof}

\begin{theorem}[Discreteness of mixing classes]
\label{thm:discrete}
Under Assumption~\ref{ass:discrete}, the set of accessible mixing
matrices
\[
  \mathcal{M}_{ST}
  = \{U_{ST} : \Xi_{ST}\ \text{admissible}\}
\]
is discrete (in fact finite) in
$\mathrm{U}(n)/(\mathrm{U}(1)^n\times\mathrm{U}(1)^n)$.
\end{theorem}

\begin{proof}
The set of admissible configurations $\Xi_{ST}$ is finite by
Assumption~\ref{ass:discrete}.
By Proposition~\ref{prop:unitary}, the map $\Xi_{ST}\mapsto U_{ST}$
is continuous, hence maps a finite set to a finite set.
Quotienting by the rephasing action (a continuous group action on a
finite set) preserves finiteness.
Thus $\mathcal{M}_{ST}$ is a finite, hence discrete, subset of
$\mathrm{U}(n)/(\mathrm{U}(1)^n\times\mathrm{U}(1)^n)$.
\end{proof}

\begin{corollary}[No continuous drift]
\label{cor:nodrift}
For any continuous deformation of boundary directions of magnitude
$\varepsilon\leq\varepsilon_0$, no continuous drift occurs in any
rephasing-invariant mixing observable.
For larger deformations that change the configuration class, the
mixing matrix can change, but only by a discrete jump to another
element of $\mathcal{M}_{ST}$.
\end{corollary}

\begin{proof}
Rephasing-invariant observables are continuous functions on
$\mathrm{U}(n)/(\mathrm{U}(1)^n\times\mathrm{U}(1)^n)$.
Their restriction to the finite set $\mathcal{M}_{ST}$ takes only
finitely many values; continuous drift is therefore impossible.
\end{proof}

\section{Connection to Hyperbolic Flavor Geometry}

The abstract framework above is concretely realized in the HFG
program~\cite{Gentry:Unified} as follows.
The group $\Gamma=\pi_1(M)$ acts on $\mathbb{H}^3$ via the
holonomy representation $\rho:\pi_1(M)\to\mathrm{PSL}(2,\mathbb{C})$.
For a loxodromic element $\gamma$, the boundary direction
$\xi(\gamma)=\hat{n}(\gamma)\in S^2$ is the geodesic axis,
extracted via the Pauli decomposition of $\log\rho(\gamma)$.
The inner product $\langle\xi(f)|\xi(g)\rangle$ becomes the
Gaussian kernel $\exp(-\theta_{fg}^2/(2\sigma^2))$, where
$\theta_{fg}=\arccos|\hat{n}(f)\cdot\hat{n}(g)|$.

For $M_{\mathrm{CKM}}=m006(-5,2)$ (OrientableClosedCensus[43],
$H_1=\mathbb{Z}/5$, vol $=2.029$), the three-word sector
$\{\mathrm{aaB},\mathrm{AbA},\mathrm{AAb}\}$ gives a mixing matrix
with Frobenius fitness $0.016482$ against PDG~2024 CKM values,
and Cabibbo angle error $0.19\%$.
The homological class collision $[\mathrm{AbA}]=[\mathrm{AAb}]=2$
in $H_1=\mathbb{Z}/5$ forces $J_{\mathrm{CP}}=0$ via
Corollary~\ref{cor:nodrift} --- consistent with the observed
suppression $|J_{\mathrm{CKM}}|\approx3\times10^{-5}$.

For $M_{\mathrm{PMNS}}=m003(-2,3)$ (OrientableClosedCensus[1],
$H_1=\mathbb{Z}/5$, vol $=0.981$), the Borel (lower-triangular)
construction on the sector $\{\mathrm{aa},\mathrm{aaB},\mathrm{baa}\}$
gives fitness $0.005087$ --- the global minimum of the construction,
with $p<10^{-4}$ against $50{,}000$ Haar-random unitary matrices.
The distinct homology classes $[\mathrm{aa}]=1$,
$[\mathrm{aaB}]=4$, $[\mathrm{baa}]=3$ (all distinct mod~5)
ensure $J_{\mathrm{CP}}\neq0$, consistent with the large leptonic
CP violation observed.

\section{Conclusions}

We have shown that the discreteness of boundary directions in a
compact hyperbolic 3-manifold implies the discreteness of all
admissible mixing matrices modulo rephasings (Theorem~\ref{thm:discrete}),
and that bounded deformations cannot produce continuous drift in
physical observables (Theorem~\ref{thm:stable}, Corollary~\ref{cor:nodrift}).
This provides a model-independent kinematic explanation for why
the CKM and PMNS mixing matrices cannot vary continuously: they
are selected from a discrete geometric set determined by the
topology of the underlying manifold.

\section*{Data Availability}
No data were generated. Companion computational scripts are at
\url{https://github.com/drmlgentry/hyperbolic-flavor-scan}.

\section*{Acknowledgments}
The author thanks the developers of SnapPy and SageMath.

\begin{thebibliography}{99}

\bibitem{Cabibbo}
N.~Cabibbo, Phys.\ Rev.\ Lett.\ \textbf{10}, 531 (1963).

\bibitem{KM}
M.~Kobayashi, T.~Maskawa,
Prog.\ Theor.\ Phys.\ \textbf{49}, 652 (1973).

\bibitem{Pontecorvo}
B.~Pontecorvo, Sov.\ Phys.\ JETP \textbf{6}, 429 (1957).

\bibitem{Maki}
Z.~Maki, M.~Nakagawa, S.~Sakata,
Prog.\ Theor.\ Phys.\ \textbf{28}, 870 (1962).

\bibitem{PDG2024}
S.~Navas et al.\ (PDG),
Phys.\ Rev.\ D \textbf{110}, 030001 (2024).

\bibitem{AltarelliFeruglio}
G.~Altarelli, F.~Feruglio,
Rev.\ Mod.\ Phys.\ \textbf{82}, 2701 (2010).

\bibitem{KingLuhn}
S.~F.~King, C.~Luhn,
Rept.\ Prog.\ Phys.\ \textbf{76}, 056201 (2013).

\bibitem{Jarlskog}
C.~Jarlskog,
Z.\ Phys.\ C \textbf{29}, 491 (1985).

\bibitem{Ratcliffe}
J.~G.~Ratcliffe,
\textit{Foundations of Hyperbolic Manifolds}, Springer (2006).

\bibitem{Gentry:CKM}
M.~L.~Gentry,
\textit{Quark Mixing from Hyperbolic Geometry},
SSRN \href{https://doi.org/10.2139/ssrn.6583550}{6583550} (2026);
submitted to J.\ Geom.\ Phys.\ (JGP13076).

\bibitem{Gentry:PMNS}
M.~L.~Gentry,
\textit{Lepton Mixing from Borel Structure of Hyperbolic Holonomy},
SSRN \href{https://doi.org/10.2139/ssrn.6583549}{6583549} (2026);
submitted to Ann.\ Henri Poincar\'e (AHPO-D-26-00255).

\bibitem{Gentry:Unified}
M.~L.~Gentry,
\textit{Hyperbolic Flavor Geometry},
SSRN \href{https://doi.org/10.2139/ssrn.6775158}{6775158} (2026);
submitted to Phys.\ Rev.\ D (DS14327).

\end{thebibliography}

\end{document}